\chapter{Introduction}%
\label{chap:intro}

This chapter serves as a guide to using this thesis template effectively. It covers the template's features, basic typographic rules, and best practices for structuring your document. Whether you're a new user or an experienced \LaTeX\ writer, these guidelines will help you produce a professional and consistent thesis.

% ---------------------------------------------------------------------------------------------------
\section{About the template}
\label{sec:about}

The iFR thesis template is designed to streamline the writing process while ensuring professional formatting and adherence to academic standards. Beyond its clean and polished appearance, this template includes comprehensive examples and built-in functionality for common thesis elements.

\subsection{Key Features}
\label{subsec:features}

This template comes with ready-to-use examples for:
\begin{itemize}
    \item How to use the nomenclature
    \item How to include images in different layouts
    \item How to cite from the bibliography (.bib) file
    \item How insert a formatted piece of code
\end{itemize}

For comprehensive information beyond this introduction, please consult the \texttt{README.md} file included with this template. This document contains essential information about:
\begin{itemize}
    \item How to build your final pdf document
    \item How to set up VSCode
\end{itemize}

\subsection{Paragraph Indentation}
\label{subsec:indentation}

\textbf{First paragraphs after headings} (chapters, sections, subsections) \textbf{should not be indented}. This is a standard typographic convention that helps visually separate section breaks from normal paragraph breaks. Subsequent paragraphs within the same section should be indented.

\subsection{Paragraph Breaks}
\label{subsec:paragraph-breaks}

\textbf{Never force line breaks manually} using \texttt{\textbackslash\textbackslash} or \texttt{\textbackslash newline} between paragraphs. Instead, separate paragraphs with a blank line in your source code. This approach:

\begin{itemize}
    \item Allows \LaTeX\ to optimize line and page breaks
    \item Ensures consistent paragraph spacing throughout the document
    \item Makes your source code more readable and maintainable
    \item Prevents formatting errors when the document is recompiled
\end{itemize}
